% =====================================================================
%  MLE-AI-201 Machine Learning -- Final Project report
%  IEEE two-column look with standard packages only (no IEEEtran.cls).
%
%  Build (after `python -m src.run_all` has written report/figures/):
%     tectonic report.tex      or      pdflatex report.tex ; pdflatex report.tex
% =====================================================================
\documentclass[10pt,twocolumn,a4paper]{article}

\usepackage[a4paper,top=0.75in,bottom=1.0in,left=0.68in,right=0.68in,columnsep=0.25in]{geometry}

\usepackage{iftex}
\ifPDFTeX                      % pdflatex (Overleaf default)
  \usepackage{mathptmx}        % Times text + math
  \usepackage[T1]{fontenc}
  \usepackage[utf8]{inputenc}
  \usepackage[safe]{tipa}      % \textschwa for the Azerbaijani letter schwa
  \DeclareUnicodeCharacter{0259}{\textschwa}
  \DeclareUnicodeCharacter{018F}{\textipa{@}}
\else                          % xelatex / lualatex / tectonic
  \usepackage{mathptmx}        % Times math
  \usepackage{fontspec}
  \setmainfont{texgyretermes}[Extension=.otf, UprightFont=*-regular,
    BoldFont=*-bold, ItalicFont=*-italic, BoldItalicFont=*-bolditalic]
                               % Times clone with all Azerbaijani letters
\fi

\usepackage{graphicx}
\usepackage{amsmath,amssymb}
\usepackage{booktabs}
\usepackage{array}
\usepackage{titlesec}
\usepackage{caption}
\usepackage{cite}
\usepackage[hidelinks]{hyperref}
\usepackage{url}

\renewcommand{\thesection}{\Roman{section}}
\renewcommand{\thesubsection}{\Alph{subsection}}
\titleformat{\section}[block]{\centering\normalsize\scshape}{\thesection.}{0.5em}{}
\titleformat{\subsection}[block]{\normalsize\itshape}{\thesubsection.}{0.5em}{}
\titlespacing*{\section}{0pt}{1.4ex plus 1ex minus .2ex}{1ex}
\titlespacing*{\subsection}{0pt}{1.0ex plus 1ex minus .2ex}{0.6ex}

\newcommand{\abstractkw}[1]{\noindent\textbf{\textit{#1}}}

\captionsetup{font=small,labelsep=period}
\setlength{\parindent}{1em}

% =====================================================================
\begin{document}

\twocolumn[
  \begin{@twocolumnfalse}
  \begin{center}
    {\LARGE\bfseries Pricing Baku Real Estate from Scratch:\\
     Decision Trees and Pegasos SVMs on bina.az Listings\par}
    \vspace{1.2em}
    {\large
      Askerova, Melek\quad Gurbanova, Sevda\quad Ibrahimova, Nezrin\par}
    \vspace{0.5em}
    {\itshape MLE-AI-201 Machine Learning, AI Academy --- Cohort II\par}
    \vspace{1.4em}
  \end{center}
  \begin{quote}
    \small
    \abstractkw{Abstract---}%
    We study two prediction tasks on 58,007 residential sale listings
    scraped from bina.az: regression of the asking price (Task~A) and
    classification of a derived price tier, premium vs.\ standard, split at
    the training median of 208,000~AZN (Task~B). We implemented a CART-style
    decision tree (Gini, entropy and MSE criteria) and a soft-margin SVM
    trained with the Pegasos sub-gradient method from scratch in NumPy. We
    also added random Fourier features (RFF) to approximate an RBF kernel.
    All models share one stratified 70/15/15 split. Hyperparameters are
    selected on validation only, and every model is compared with its
    scikit-learn counterpart under matched settings. Our tree reaches test
    $R^{2}=0.709$ on price (sklearn: $0.710$) and F1 $=0.926$ on the tier
    task (sklearn: $0.925$). Linear Pegasos reaches the same objective value
    as \texttt{LinearSVC} (F1 $0.905$ vs.\ $0.906$). RFF Pegasos is the best
    SVM (F1 $0.912$, AUC $0.974$) and matches an exact RBF \texttt{SVC}.
    Trees win on F1 because they capture area$\times$district interactions;
    the SVMs rank listings better (higher AUC). Errors concentrate on luxury
    listings, which the tree under-prices by about 38\%, and on rare
    districts. One command reproduces every number:
    \href{https://github.com/nezrinibrahimova/ML_Final_Project}{\texttt{github.com/nezrinibrahimova/ML\_Final\_Project}}
    (tag \texttt{v1.0-final}).

    \vspace{0.6em}
    \abstractkw{Index Terms---}decision tree, support-vector machine,
    Pegasos, random Fourier features, regression, classification,
    real-estate price prediction.
  \end{quote}
  \vspace{1.0em}
  \end{@twocolumnfalse}
]

% =====================================================================
\section{Introduction}
Asking prices on Baku's largest listing site, bina.az, vary by more than
two orders of magnitude, from 16.5k to 5.95M~AZN. Size, building type
and, above all, location drive them. A model that predicts the price, or
says whether a listing belongs to the expensive half of the market, is
useful to buyers and sellers. It is also a good test bed for two classic
learners whose inductive biases differ: axis-aligned trees and
maximum-margin linear classifiers.

Our contributions are:
(i)~a from-scratch NumPy decision tree for both tasks, with a vectorised
split search and four overfitting controls;
(ii)~a from-scratch Pegasos SVM with a stable learning-rate offset, margin
and support-vector reporting, and an RFF kernel approximation;
(iii)~a single evaluation protocol shared by every model, with matched
scikit-learn baselines, 5-fold CV and timing;
(iv)~an error analysis by price segment, district and market cluster;
(v)~two bonuses: closed-form Ridge regression, and PCA with $k$-means
market segmentation.

\section{Background}
A decision tree recursively splits the feature space with axis-aligned
tests $x_j\le\theta$. At every node it chooses the test that maximises
the impurity decrease
$\Delta I=I(P)-\tfrac{N_L}{N}I(L)-\tfrac{N_R}{N}I(R)$. Here $I$ is the Gini
index $1-\sum_k p_k^2$, the entropy $-\sum_k p_k\log_2p_k$ (its decrease
is the information gain) or, for regression, the variance (MSE)~\cite{cart}.
Trees need no feature scaling and model interactions naturally, but deep
trees have high variance.

The soft-margin SVM finds the hyperplane with the widest margin
$2/\lVert\mathbf w\rVert$ while paying a hinge loss for points inside the
margin. Pegasos~\cite{pegasos} solves its primal by stochastic
sub-gradient descent, which avoids a QP solver. Random Fourier
features~\cite{rff} approximate a shift-invariant kernel with an explicit
finite-dimensional map, so a linear SVM can learn an approximately
RBF-kernel decision boundary.

% =====================================================================
\section{Data and Preprocessing}
\emph{Source.} The Kaggle \emph{binaaz-sale-project} dump contains
100,775 rows and 51 columns, scraped between 5~Oct and 19~Nov~2024. All
prices are in AZN, so no currency conversion is needed.

\emph{Cleaning} (\texttt{src/data\_prep.py}). (1)~We keep the three
residential categories (new build, old build, house/villa; 91,567 rows)
and drop land, commercial, office and garage listings. (2)~Listings
appear repeatedly across scrapes, so we keep the latest record per
listing URL (58,452 rows). (3)~We parse area, floor/total floors and land
size from their text fields. (4)~Outlier rules: area $15$--$1000$~m$^2$,
price per m$^2$ $300$--$10{,}000$~AZN (computed only for filtering and
then discarded), land $\le 1000$ sot, valid room count, and coordinates
inside Azerbaijan. (5)~Missing floor and land values become 0 (houses
have no floor; flats have no land). This leaves 58,007 listings.

\emph{Leakage.} \texttt{unit\_price} (price per m$^2$) and
\texttt{total\_price} are functions of the target and are dropped. Free
text, owner and shop names, raw address, view counts and image URLs are
not used as features either.

\emph{Features.} Ten numeric features (area, rooms, floor, total floors,
land, latitude, longitude, and repair/deed/mortgage flags) plus one-hot
category and district. The one-hot columns are defined on the training
split only (unseen levels map to all-zeros), giving 148 features.

\emph{Targets and split.} Prices are right-skewed (median 208,000, mean
263,374, max 5.95M~AZN), so the regression tree is trained on
$\log(\text{price})$ and its metrics are reported in AZN. The tier label
is premium if the price exceeds the \emph{training} median (208,000~AZN).
One stratified split on price quartiles (seed 42) is shared by every
model (Table~\ref{tab:data}). A unit test asserts that the tree and SVM
scripts use identical rows.

\begin{table}[t]
  \centering
  \caption{Dataset summary after cleaning.}
  \label{tab:data}
  \begin{tabular}{@{}lr@{}}
    \toprule
    Property & Value \\
    \midrule
    Raw rows / columns          & 100,775 / 51 \\
    Rows after cleaning         & 58,007 \\
    Features used               & 148 (10 numeric + one-hot) \\
    Districts                   & 138 \\
    Train / Val / Test          & 40,605 / 8,701 / 8,701 \\
    Tier threshold (train median) & 208,000 AZN \\
    Premium share (train/val/test) & 50.0 / 50.0 / 50.0\,\% \\
    \bottomrule
  \end{tabular}
\end{table}

% =====================================================================
\section{Method}
\subsection{Decision tree (from scratch)}
\texttt{DecisionTree} (\texttt{src/decision\_tree.py}) supports
classification with Gini or entropy and regression with MSE. Growth is
recursive. A node becomes a leaf when it is pure, when
\texttt{max\_depth} is reached, when it has fewer than
\texttt{min\_samples\_split} samples, or when the best split's weighted
impurity decrease $\tfrac{N}{N_{\text{tot}}}\Delta I$ is below
\texttt{min\_impurity\_decrease}. \texttt{min\_samples\_leaf} masks
candidate splits that would create a smaller child. A unit test checks
that each of the four hyperparameters changes the fitted tree.

\emph{Split search.} For every feature we sort once and compute
cumulative class counts (classification) or cumulative $\sum y$ and
$\sum y^2$ (regression). Every candidate threshold is then scored in
$O(1)$ and all thresholds at once, in vectorised form, so a node costs
$O(d\,n\log n)$. Thresholds are midpoints between distinct consecutive
values, as in scikit-learn. A guard handles adjacent floats whose midpoint
rounds up to the larger value, which would otherwise create an empty
child; a regression test covers this case.

\emph{Outputs.} Leaves store the class frequencies, which give
\texttt{predict\_proba}, or the mean target. \texttt{feature\_importances\_}
sums each feature's weighted impurity decreases and normalises them. A
unit test checks that our root split, predictions and importances match
scikit-learn on small deterministic data.

\subsection{Support-vector machine (from scratch)}
Our \texttt{PegasosSVM} (\texttt{src/svm.py}, NumPy only) minimises the
soft-margin primal with labels $y_i\in\{-1,+1\}$ (premium $=+1$):
\begin{equation}
  J(\mathbf{w},b)=\tfrac{\lambda}{2}\lVert\mathbf{w}\rVert^{2}
  +\tfrac{1}{n}\sum_{i=1}^{n}\max\!\big(0,\,1-y_i(\mathbf{w}^{\top}\mathbf{x}_i+b)\big).
  \label{eq:svm}
\end{equation}
At step $t$ we draw a mini-batch $A_t$ of size $k$ and let
$A_t^{+}=\{i\in A_t: y_i(\mathbf{w}^{\top}\mathbf{x}_i+b)<1\}$ be the
points with active hinge. The sub-gradient update is
\begin{align}
  \mathbf{w} &\leftarrow (1-\eta_t\lambda)\,\mathbf{w}
     + \tfrac{\eta_t}{k}\textstyle\sum_{i\in A_t^{+}} y_i\mathbf{x}_i,\\
  b &\leftarrow b + \tfrac{\eta_t}{k}\textstyle\sum_{i\in A_t^{+}} y_i .
\end{align}
The bias is not regularised, so it receives a plain sub-gradient step.
Each epoch is one shuffled pass over the training set; every few steps
we record $J(\mathbf{w},b)$ on the full training set to show that it
decreases.

\emph{Learning rate.} The paper's schedule
$\eta_t=1/(\lambda t)$~\cite{pegasos} makes the first step
$\eta_1=1/\lambda$ (e.g.\ $10^{3}$ for $\lambda=10^{-3}$). The factor
$(1-\eta_t\lambda)$ pulls $\mathbf{w}$ back after such a step, but
nothing pulls the unregularised bias back. On separable toy data one
early violator moved $b$ to $\approx 370$, and it was still there after
10 epochs. We therefore use $\eta_t=1/(\lambda(t+t_0))$ with
$t_0=1/\lambda$, so that $\eta_1\approx1$. We also implemented
$t_0=0$, $\eta_0/\sqrt t$ and a constant rate, and compare them in
Fig.~\ref{fig:svm-obj}.

\emph{Margin and support vectors.} After training we report
$\lVert\mathbf{w}\rVert$, the geometric margin $2/\lVert\mathbf{w}\rVert$
and the support vectors, i.e.\ the training points with
$y_i(\mathbf{w}^{\top}\mathbf{x}_i+b)\le 1$ (on or inside the margin, or
misclassified).

\emph{Non-linear extension: random Fourier features.} Following
Rahimi and Recht~\cite{rff}, we map
$\phi(\mathbf{x})=\sqrt{2/D}\cos(W^{\top}\mathbf{x}+\mathbf{c})$ with
$W_{jk}\sim\mathcal N(0,2\gamma)$ and $c_k\sim U[0,2\pi]$, so that
$\phi(\mathbf{x})^{\top}\phi(\mathbf{x}')\approx
\exp(-\gamma\lVert\mathbf{x}-\mathbf{x}'\rVert^{2})$. Linear Pegasos on
$\phi(X)$ then approximates an RBF-kernel SVM; a unit test checks the
kernel approximation error ($<0.05$ at $D=20{,}000$). We chose RFF over
kernelised Pegasos because one epoch costs $O(nD)$ instead of $O(n^{2})$
kernel evaluations. With $n\approx 4\times10^{4}$ training rows this
difference matters, and prediction does not need the support vectors.

\emph{Feature scaling.} Pegasos is not scale-invariant, so features are
standardised with training statistics only. For the linear SVM we
standardise every column. For the RBF models (RFF and the sklearn
\texttt{SVC} baseline) we standardise only the numeric columns and keep
the one-hot district/category columns as $0/1$. A standardised rare
district dummy becomes $\approx 200$ for the few listings in that
district, and squared distances are then dominated by district
identity. Section~\ref{sec:svm-results} measures this effect.

% =====================================================================
\section{Experimental Setup}
\emph{Metrics} are implemented in NumPy (\texttt{src/evaluate.py}):
RMSE, MAE and $R^2$ in AZN for Task~A (plus RMSE and $R^2$ on the log
scale), and precision, recall and F1 for the premium class, plus
rank-based ROC-AUC, for Task~B. AUC uses leaf frequencies for trees and
signed margins for SVMs.

\emph{Protocol.} Every hyperparameter is chosen on the validation split
and the test split is scored once, for the final configurations only.
Trees: depth and leaf sweeps for each criterion
(Fig.~\ref{fig:tree-depth}), then a joint search over
$\texttt{max\_depth}\in\{10,12,15,20\}\times
\texttt{min\_samples\_leaf}\in\{1,3,5,10\}$. SVMs: a $\lambda$ sweep, then
a $\gamma\times\lambda$ sweep for RFF and a sweep over the number of
random features $D$.

\emph{Baselines.} scikit-learn \texttt{DecisionTreeRegressor/Classifier}
use the same selected depth, leaf size and criterion. \texttt{LinearSVC}
(hinge) and \texttt{SGDClassifier} (hinge, $\alpha=\lambda$) use
$C=1/(\lambda n)$ so that the objective equals Eq.~(\ref{eq:svm}).
\texttt{SVC} (RBF) uses the selected $\gamma$ and is fitted on a
stratified 10k-row subsample, because exact kernel training is
$O(n^{2})$. Each model is also evaluated with 5-fold CV on the training
split; scaling is fitted inside every fold (an sklearn \texttt{Pipeline}
for the baselines). We report fit and predict times on one laptop CPU.

\emph{Reproducibility.} Every source of randomness uses seed 42.
\texttt{python -m src.run\_all} rebuilds every number, table and figure
from a fresh clone with the pinned \texttt{requirements.txt} in about
24~minutes; two independent runs produced identical tables.

% =====================================================================
\section{Results}
Table~\ref{tab:results} is the headline comparison on the test split.

\begin{table}[t]
  \centering
  \caption{Headline test results (8,701 listings). Task~A in AZN;
  Task~B for the premium class. CV: 5-fold on the training split.}
  \label{tab:results}
  \footnotesize
  \setlength{\tabcolsep}{3pt}
  \begin{tabular}{@{}lccc@{}}
    \toprule
    \multicolumn{4}{@{}l}{\emph{Task A: price regression}} \\
    Model & RMSE & MAE & $R^2$ \\
    \midrule
    Train-mean baseline              & 222,308 & 131,332 & $0.000$ \\
    Ours: Ridge (bonus)              & 124,377 & 57,792  & $0.687$ \\
    Ours: decision tree (log target) & 119,931 & 43,889  & $0.709$ \\
    sklearn tree                     & 119,699 & 43,852 & $0.710$ \\
    \midrule
    \multicolumn{4}{@{}l}{\emph{Task B: price tier}} \\
    Model & F1 & AUC & CV F1 \\
    \midrule
    Ours: tree, Gini        & $.923$ & $.954$ & $.920\pm.004$ \\
    Ours: tree, entropy     & $\mathbf{.926}$ & $.930$ & $.920\pm.003$ \\
    sklearn tree, Gini      & $.921$ & $.953$ & $.920\pm.004$ \\
    sklearn tree, entropy   & $.925$ & $.929$ & $.920\pm.004$ \\
    Ours: Pegasos linear    & $.905$ & $.964$ & $.904\pm.005$ \\
    Ours: Pegasos + RFF     & $.912$ & $\mathbf{.974}$ & n/a \\
    sklearn \texttt{LinearSVC} & $.906$ & $.964$ & $.905\pm.005$ \\
    sklearn \texttt{SGDClassifier} & $.906$ & $.964$ & $.905\pm.005$ \\
    sklearn \texttt{SVC} RBF (10k) & $.912$ & $.973$ & n/a \\
    \bottomrule
  \end{tabular}
\end{table}

\subsection{Decision tree}
\emph{Hyperparameters.} Fig.~\ref{fig:tree-depth} shows the classic
bias--variance picture. Training F1 climbs to $0.99$ at depth 20, while
validation F1 rises steeply up to depth 6 and then flattens
($0.920$--$0.923$ at depth 12, $0.925$--$0.927$ at 20). Gini and entropy
are within $0.007$ of each other at every depth. The joint search selected depth 20
for all three trees, with leaf sizes 10 (regression), 5 (Gini) and 1
(entropy); validation scores were RMSE$_{\log}=0.223$, F1 $0.927$ and
$0.927$. Depth 20 is the edge of our grid, but the validation gain from
15 to 20 is only about $0.002$ F1, so a deeper grid would change little.

\emph{Comparison with scikit-learn.} With identical settings our trees
reproduce scikit-learn closely. Leaf counts are 2,841 vs.\ 2,837
(regression), 1,506 vs.\ 1,511 (Gini) and 2,326 vs.\ 2,350 (entropy).
Test predictions agree on 99.5\% (Gini) and 98.3\% (entropy) of listings,
and the regression predictions correlate at $0.9995$. 9--10 of the top-10
important features coincide; area dominates (importance $0.62$),
followed by longitude and latitude. The small differences come from ties:
when two splits have exactly the same impurity decrease, scikit-learn
breaks the tie with a random feature permutation, while we keep the first.
CV confirms the match (Table~\ref{tab:results}). Our tree is 7--22$\times$
slower to fit (15.3\,s vs.\ 2.1\,s on the full training set), the price of
Python-level recursion compared with Cython.

\emph{Gini vs.\ entropy.} The entropy tree has the best F1 ($0.926$) but
a much lower AUC ($0.930$ vs.\ $0.954$). With
\texttt{min\_samples\_leaf}$=1$ its leaves are pure, so almost every
probability is 0 or 1, which gives a poor ranking. The Gini tree (leaf
size 5) is the better-calibrated choice at the same F1.

\begin{figure}[t]
  \centering
  \includegraphics[width=\linewidth]{figures/tree_B_depth.png}
  \caption{Task~B tree: train (dashed) and validation (solid) F1 against
  \texttt{max\_depth} for Gini and entropy.}
  \label{fig:tree-depth}
\end{figure}

\subsection{SVM (Task B)}\label{sec:svm-results}
\emph{Protocol.} We sweep $\lambda\in\{10^{-6},\dots,10^{-1}\}$ for linear
Pegasos (100 epochs, $k=256$). For RFF Pegasos (30 epochs, $k=64$) we sweep
$\gamma\in\{0.01,0.03,0.1,0.3\}\times\lambda\in\{10^{-6},10^{-5},10^{-4}\}$
at $D=1000$, then vary $D$; the final model uses $D=2000$. All choices
maximise validation F1.

\emph{The from-scratch solver reaches the optimum.} At the selected
$\lambda=10^{-3}$, Pegasos ends at $J=0.2619$, which equals the objective
of the \texttt{LinearSVC} solution. Test F1 differs by $0.001$, and the CV
F1 difference is $0.0008$, well inside one standard deviation
(Table~\ref{tab:results}). Fig.~\ref{fig:svm-obj} shows $J$ decreasing.
With the paper's schedule ($t_0=0$), $J$ first jumps above its starting
value of 1 before recovering, while our offset schedule drops at once.
$\eta_0/\sqrt t$ is too slow at the same budget ($J=0.291$, val F1
$.889$).

\emph{Margin trade-off.} Increasing $\lambda$ widens the margin
$2/\lVert\mathbf{w}\rVert$ from $0.33$ ($\lambda=10^{-6}$) to $1.69$
($\lambda=10^{-1}$). The number of support vectors grows from $9.1$k to
$19.8$k (22\% $\to$ 49\% of the training set). Validation F1 peaks at
$\lambda=10^{-3}$ (margin $0.48$, F1 $.908$) and falls to $.862$ at the
widest margin, where the model underfits (Fig.~\ref{fig:svm-lambda}).

\emph{Non-linearity helps.} RFF Pegasos is the best SVM on test
(F1 $.912$, AUC $.974$). It matches the exact RBF \texttt{SVC} and keeps
training linear in $n$. Validation F1 rises with the number of random
features: $.888$ at $D=250$, $.901$ at $500$, $.913$ at $1000$ and $.917$
at $2000$. The approximation is the bottleneck at small $D$. The selected
bandwidth is $\gamma=0.1$; a larger $\gamma=0.3$ overfits to local
neighbourhoods (val F1 $.890$). Feature scaling matters as much as the
kernel. With every column standardised (largest $|x|=202$, a rare district
dummy), the same RFF model reaches only val F1 $.859$ / AUC $.933$. With
one-hot columns kept as $0/1$ it reaches $.913$ / $.973$.

\emph{What the linear SVM uses.} The weights of the standardised features
are dominated by \texttt{area\_m2} ($w=3.7$). The next largest weights are
district effects: outer settlements (Masazır, Binəqədi, Binə) push toward
\emph{standard}, central districts (Nizami, 28 May, Ağ şəhər) toward
\emph{premium}, and new-build vs.\ house/villa category matters as well.

\begin{figure}[t]
  \centering
  \includegraphics[width=\linewidth]{figures/svm_objective.png}
  \caption{Training objective $J(\mathbf{w},b)$ of linear Pegasos
  ($\lambda=10^{-3}$) under four learning-rate schedules.}
  \label{fig:svm-obj}
\end{figure}

\begin{figure}[t]
  \centering
  \includegraphics[width=\linewidth]{figures/svm_lambda_sweep.png}
  \caption{Linear Pegasos: validation F1/AUC (left axis) and margin
  $2/\lVert\mathbf{w}\rVert$ (right axis, log) against $\lambda$.}
  \label{fig:svm-lambda}
\end{figure}

\subsection{Error analysis}
\emph{Luxury listings.} The regression tree's median absolute percentage
error is 8.0\% for the cheaper half of the test set and 8.3\% for the
50--90th percentiles. It rises to 15.2\% for the 90--99th percentiles and
34.7\% for the top 1\% (mean price 1.7M~AZN). There the mean log error is
$-0.48$, so the tree under-prices luxury listings by about 38\%
(Fig.~\ref{fig:errors}, left). A leaf averages many listings, so the few
extreme prices are pulled toward their neighbours.

\emph{Rare districts.} The districts with the largest errors have few
training listings. For example, Xaçmaz (21 training listings) has a median
error of 122\%, and Qala and Kürdəxanı (20--34 listings) have about
60\%. In such districts the tree must fall back on area and coordinates
learned elsewhere.

\emph{The tier border.} Task~B errors are mostly near the threshold. The
error rate is 27.1\% for listings within 10\% of 208,000~AZN, 12.8\% at
10--25\%, and below 4\% beyond 25\% (Fig.~\ref{fig:errors}, right). Most
of the remaining classification error is therefore the irreducible
ambiguity of an arbitrary cut on a noisy price, not a modelling failure.

\begin{figure}[t]
  \centering
  \includegraphics[width=\linewidth]{figures/tree_error_analysis.png}
  \caption{Tree error analysis on test. Left: Task~A median percentage
  error by true-price segment. Right: Task~B error rate against distance
  from the tier threshold.}
  \label{fig:errors}
\end{figure}

\subsection{Bonus: Ridge regression and market segments}
\emph{Ridge (P8).} Closed-form Ridge,
$\hat{\mathbf w}=(X^{\top}X+\alpha I)^{-1}X^{\top}\mathbf y$ on standardised
features ($\alpha=0.01$ on validation), reaches test RMSE 124,377 and
$R^2=0.687$. It matches sklearn \texttt{Ridge} to $2\times10^{-5}$~AZN, but
it is clearly behind the tree in MAE (57.8k vs.\ 43.9k). A linear price
model cannot express that one m$^2$ is worth more in the centre.

\emph{Segments (P10).} PCA keeps 125 components (90.4\% of the variance),
and $k$-means ($k=5$) on them yields market segments. The hardest segment
is the large-house cluster (6,404 training listings, median 180~m$^2$):
validation $R^2$ is $0.56$ for both Ridge and the tree, against
$0.65$--$0.93$ in the other clusters. The tree has the lower MAE in four of
five clusters. For the tier task the linear SVM has the higher AUC in
every cluster, consistent with the ranking advantage in
Table~\ref{tab:results}. These cluster numbers are diagnostics on the
validation split, not independent test estimates.

% =====================================================================
\section{Discussion and Limitations}
\emph{Trees vs.\ SVMs.} The tree wins on F1 ($0.923$--$0.926$
vs.\ $0.905$--$0.912$). The price tier is largely an interaction of size
and location, which axis-aligned splits express directly: within one
district, area decides. The SVMs win on AUC ($0.964$--$0.974$
vs.\ $0.930$--$0.954$), because a continuous margin ranks listings
smoothly while a tree has a few thousand discrete leaf scores. A user who
needs a ranking or a calibrated score should prefer the RFF SVM; a user
who needs the hard premium/standard label should prefer the tree.

\emph{Bias and variance.} The deep trees (depth 20) have training F1
near $0.99$ but validation F1 near $0.925$, so they have high variance.
The Task~A learning curve is still improving at 40k listings (validation
RMSE$_{\log}$ falls from $0.245$ to $0.240$ between 28k and 40k), while Task~B has
flattened, so averaging trees (a forest) is the more promising remedy.
Linear Pegasos is the opposite case: its CV, validation and test F1 all
lie within $0.005$ ($0.904$--$0.908$), so its errors are mostly bias.
RFF trades some of that bias for variance and gains $0.007$ F1.

\emph{Limitations.} (i)~Asking prices are not transaction prices.
(ii)~The data covers six weeks of 2024, so we cannot measure temporal
drift. (iii)~Districts with fewer than ${\sim}50$ listings are poorly
served by every model. (iv)~The exact RBF \texttt{SVC} baseline saw only
10k training rows, so the comparison with RFF (40k rows) favours ours.
(v)~Our tree is about 7$\times$ slower than scikit-learn. (vi)~The tree
grid stops at depth 20.

\section{Conclusion}
From-scratch NumPy implementations of a CART tree and a Pegasos SVM
match scikit-learn under one leakage-free protocol: tree $R^2$
$0.709$ vs.\ $0.710$, and the same SVM objective value. On bina.az
listings, trees are the most accurate classifiers of the price tier
(F1 $0.926$), while an RFF-approximated RBF SVM gives the best ranking
(AUC $0.974$). Three choices mattered as much as the model: a stable
Pegasos step size, keeping one-hot features unscaled for RBF kernels, and
a log target for prices. Next we would add a random forest of our own
trees to cut the tree's variance, add district-level target statistics
learned on training data only for rare districts, and model luxury
listings separately.

\section*{Tools and Acknowledgements}
Data: the Kaggle dataset \emph{binaaz-sale-project} by
\texttt{sehriyarmemmedli}. Libraries: NumPy, pandas, Matplotlib, and
scikit-learn, used only for baselines, CV splitting and the bonus
comparisons. AI assistance: the SVM part (\texttt{src/svm.py},
\texttt{src/svm\_experiment.py}, \texttt{tests/test\_svm.py}), the
\texttt{run\_all} integration and the drafting of this report were done
with the help of Claude Code (Anthropic). It was used to write and debug
code, run the experiments and draft text; every number in this report
was produced by \texttt{run\_all} and checked by the team.
% TODO (A, B): add one sentence each on whether and how you used AI tools
% for your parts (data/EDA/metrics/bonuses; decision tree), then delete this comment.

% =====================================================================
\begin{thebibliography}{9}
\bibitem{cart}
L.~Breiman, J.~Friedman, R.~Olshen, and C.~Stone,
\emph{Classification and Regression Trees}.
Wadsworth, 1984.
\bibitem{pegasos}
S.~Shalev-Shwartz, Y.~Singer, N.~Srebro, and A.~Cotter,
``Pegasos: primal estimated sub-gradient solver for SVM,''
\emph{Mathematical Programming}, vol.~127, no.~1, pp.~3--30, 2011.
\bibitem{rff}
A.~Rahimi and B.~Recht,
``Random features for large-scale kernel machines,''
in \emph{Advances in Neural Information Processing Systems}, vol.~20, 2007.
\end{thebibliography}

\end{document}
